\documentclass[11pt]{article}
\usepackage[T1]{fontenc}
\usepackage{lmodern,amsmath,amssymb,amsthm,geometry,microtype,xcolor}
\geometry{margin=1in}
\usepackage[colorlinks=true,allcolors=blue!45!black]{hyperref}
\setlength{\emergencystretch}{3em}
\newtheorem{theorem}{Theorem}
\newcommand{\mex}{\operatorname{mex}}
\title{A direct Lean certificate for Rust's\newline game-sequence counterexample at index twenty}
\author{Myeongjin Kang}
\date{October 26, 2026}
\begin{document}\maketitle
\begin{abstract}
Rust conjectured an identity involving two greedily constructed game
sequences. Shallit disproved it in 2022 and identified index twenty as a
counterexample. We verify that same finite witness directly in Lean,
constructing the complete infinite sequences and proving their exact
greedy selection rule. Their certified values are $X(20)=30$, $Y(20)=56$,
$X(56)=86$ and $Y(30)=84$. The proposed equality therefore reads $170=171$.
The certificate uses exact integer computation, rather than an assumed
automaton or a table detached from the source definition.
\end{abstract}
\noindent\textit{Credit and scope.} This is verification of Shallit's known
counterexample, not a new disproof or first-formalization claim. The source
also establishes infinitely many counterexamples; that stronger statement
is not claimed here. This AI-assisted paper has original prose CC BY 4.0
and original code MIT.

\section{The actual recursive sequences}
Definition 1 of Shallit's \emph{Some Tribonacci Conjectures} starts with
$(X(0),Y(0))=(0,0)$ and $(X(1),Y(1))=(1,2)$~\cite{Shallit}.
At each subsequent index $n$, select positive integers $x<y$ such that
neither appears among the earlier $X$ or $Y$ values, and $y-x$ appears
among neither the earlier differences nor earlier sums. Choose $x$ as
small as possible, then $y$ as small as possible for that $x$. This
lexicographic reading is explicitly used in the source's Section 3.2
verification, rather than being an additional convention imposed here.

We give a total construction with finite sets. Immediately before choosing
pair $n+1$, let
\[
 U_n=\{X(k),Y(k):0\le k\le n\},\qquad
 D_n=\{Y(k)-X(k),Y(k)+X(k):0\le k\le n\}.
\]
Both sets contain zero. With $\mex(S)$ denoting the least nonnegative
integer outside a finite set $S$, put
\begin{align*}
 x&=\mex(U_n),\\
 E_n&=D_n\cup\{\max(u-x,0):u\in U_n\},\\
 d&=\mex(E_n),\qquad (X(n+1),Y(n+1))=(x,x+d).
\end{align*}
Update $U_n$ by inserting $x,x+d$ and $D_n$ by inserting $d,2x+d$.
The initial state is $U_0=D_0=\{0\}$.

Every mex exists: if $S$ is finite, one plus its largest element is
missing (with largest value zero for an empty set). In Lean the mex is
computed by taking the minimum of the missing elements in
$\{0,\ldots,\max S+1\}$. We prove both that this number is missing and
that every smaller nonnegative number is present. Thus the construction
is defined at every index, not just through the finite witness.

\section{Why this is precisely the source construction}
Since $0\in U_n\cap D_n$, both $x$ and $d$ above are positive.
The choice gives $x\notin U_n$ and $d\notin D_n$. Also $x+d\notin U_n$:
otherwise $d=(x+d)-x$ would belong to $E_n$. Hence $(x,x+d)$ is admissible.
Any other admissible pair $(u,v)$ has $u\notin U_n$, so $x\le u$.
If $u=x$, then $v-x>0$ lies outside $D_n$ and outside
$\{\max(w-x,0):w\in U_n\}$: an equality with a truncated difference,
being positive, would force $w=v$, contrary to $v\notin U_n$.
Therefore $d\le v-x$ and $x+d\le v$. This proves the exact two-stage
minimality rule.

Induction on the state updates proves the displayed history-set formulas.
The first choice is $(1,2)$, as required. Including history index zero
adds only zero to each forbidden set; all tested numbers $x,y,y-x$ are
positive, so this changes no admissible pair. The formal proof explicitly
identifies these histories with zero inserted into the positive-index
histories of Definition 1, and proves that admissibility and least-pair
selection are unchanged by this insertion.

Least admissible pairs are unique: applying each pair's minimality to the
other first equates their first coordinates, then their second coordinates.
Strong induction on the history consequently proves uniqueness of the
entire sequence satisfying the source selection rule. The construction
above therefore computes the actual recursively defined sequences.

\section{The exact finite witness}
Here are the successive pairs through index 56; within a row they are
listed in increasing index order. The omitted index-zero pair is $(0,0)$.
Every entry is obtained by the finite mex rule just proved.
\[
\renewcommand{\arraystretch}{1.2}
\begin{array}{c|rrrrrrrr}
 n&\multicolumn{8}{c}{(X(n),Y(n))}\\\hline
1\text{--}8 & (1,2) & (3,5) & (4,8) & (6,11) & (7,13) & (9,16) & (10,19) & (12,22) \\
9\text{--}16 & (14,25) & (15,28) & (17,31) & (18,33) & (20,36) & (21,39) & (23,42) & (24,45) \\
17\text{--}24 & (26,48) & (27,50) & (29,53) & (30,56) & (32,59) & (34,62) & (35,65) & (37,68) \\
25\text{--}32 & (38,70) & (40,73) & (41,76) & (43,79) & (44,81) & (46,84) & (47,87) & (49,90) \\
33\text{--}40 & (51,93) & (52,96) & (54,99) & (55,101) & (57,104) & (58,107) & (60,110) & (61,113) \\
41\text{--}48 & (63,116) & (64,118) & (66,121) & (67,124) & (69,127) & (71,130) & (72,133) & (74,136) \\
49\text{--}56 & (75,138) & (77,141) & (78,144) & (80,147) & (82,150) & (83,153) & (85,156) & (86,158) \\
\end{array}
\]
The finite table can be verified by scanning unused positive integers and
unused differences at each step. The four required witness entries are also checked in Lean's
kernel on the actual recursive definitions; no table values are axioms or
hypotheses.

\begin{theorem}
The conjectured identity
\[
 X(Y(n))+Y(X(n))=2X(n)+2Y(n)-1\qquad(n\ge1)
\]
fails at $n=20$.
\end{theorem}
\begin{proof}
The certified entries give
\[
 X(20)=30,\quad Y(20)=56,\quad X(56)=86,\quad Y(30)=84.
\]
The two sides are $86+84=170$ and
$2\cdot30+2\cdot56-1=171$, respectively, so they are unequal.
\end{proof}

\section{Formal coverage and reproducibility}
The module \texttt{Rust.lean} defines a finite computable mex, the state
updates, and the complete functions $X,Y$ on natural numbers. It proves
mex minimality, source admissibility, least-pair uniqueness, the actual
history formulas, the source rule, global sequence uniqueness, and the
original positive-history/seed bridge. The exact witness declaration
\texttt{exact\_values} uses kernel-checked decidable integer computation.
The main declaration \texttt{rust\_counterexample} combines source-rule
validity, all four seeds, admissibility of index twenty, and failure of
the proposed equality. There are no imported witness values or automaton
correctness assumptions.

Lean 4.34.1 and the pinned mathlib dependencies are required. The isolated
verifier compiles all source bytes and audits every local declaration's
type and transitive axioms. Source provenance, independent review and
reproduction instructions accompany the paper. MathDB 356369 records
this already refuted conjecture~\cite{MathDB}; only one source question
is counted, rather than one outcome per computed entry.
\begin{thebibliography}{9}
\bibitem{Shallit} J. Shallit, \emph{Some Tribonacci Conjectures},
arXiv:2210.03996v1, October 8, 2022, Definition 1, Section 3.2 and
Section 4, Conjecture 5 with its disproof.
\url{https://arxiv.org/abs/2210.03996v1}.
\bibitem{MathDB} MathDB, \emph{Rust's conjecture on the Tribonacci
combinatorial-game sequences}, problem 356369, accessed October 26, 2026.
\url{https://mathdb.com/p/356369}.
\end{thebibliography}
\end{document}
